\section{User Interface and Acceptance Planning}

\subsection{Interface Design}

The interface is organised around a single principle: the operator should never be
required to scrub. Every screen terminates in a moment that plays.
Figure~\ref{fig:ui} shows the four principal states of the application --- the project workspace, the ingestion
dialog, the agent chat with its clip reel and timeline studio, and the per-recording
detail view.

\optfigure{falcon-vqa-ui-design.png}{Operator interface. Clockwise from top left: the
project workspace and video grid; the ingestion dialog, whose analyzer list is built
from the analysis service's own capability endpoint; the per-recording detail view with
its chunk map and tabbed record; and the agent chat with the clip reel artifact panel
and the timeline studio.}{1.0}{fig:ui}

\begin{table}[htbp]
\centering
\small
\begin{tabularx}{\textwidth}{@{}L{3.4cm}X@{}}
\toprule
\textbf{Screen} & \textbf{Operator activity} \\
\midrule
Landing and documentation & Presents the system and hosts the full documentation: concepts, workspace guides, and the API reference. \\
Authentication  & Registration and sign-in. Every subsequent route is guarded, and every query is restricted to the signed-in account. \\
Projects        & Lists projects, each grouping a set of recordings. Opening one enters its workspace. \\
Workspace       & Adds recordings by file or by link, selects the analyzers and the chunking mode, and observes each recording advance through \emph{Uploading}, \emph{Indexing} and \emph{Ready}. Only ready recordings may be searched, and the interface states this rather than returning an empty result. \\
Agent chat      & Accepts a description or a question, and a selection of recordings to consider. The assistant chooses the retrieval operations, and answers remain concise and cite timecodes. Conversations persist and may be resumed. \\
Clip reel panel & Opens beside the chat whenever the results are moments. A player is placed above the list of moments; selecting one plays it, and they may be played consecutively as a reel. \\
Timeline studio & Presents the recording's scenes, transcript and chapters against a timeline that follows the playhead, for reviewing what was indexed in the region being watched. \\
Recording detail & Five tabs over one recording: \emph{Overview} for the video-level output, \emph{Chunks} for the chunk map and per-chunk records, \emph{Speech} for the transcript, \emph{People \& objects} for linked entities and their timelines, and \emph{Record} for the raw stored document. \\
\bottomrule
\end{tabularx}
\end{table}

\subsection{Acceptance Test Plan}

Each test is executed against a live server on the datasets listed in the following
subsection, and traces to the requirement it covers.

\begin{table}[htbp]
\centering
\small
\begin{tabularx}{\textwidth}{@{}lL{2.9cm}XL{1.6cm}@{}}
\toprule
\textbf{ID} & \textbf{Test} & \textbf{Passes when} & \textbf{Covers} \\
\midrule
AT-1 & Account isolation & A second account signing in sees none of the first account's projects, recordings or conversations, and cannot reach them by identifier. & FR-1, NFR-9 \\
AT-2 & Add a recording & A file and a link both reach \emph{Ready}, and the analyzers that ran are those selected. & FR-2 \\
AT-3 & Progress is visible & Stage and progress update during processing, and the interface remains usable throughout. & FR-3, NFR-3 \\
AT-4 & Search by description & A described event returns the correct moment within the leading results. & FR-4 \\
AT-5 & Filtered search & Restricting by time and by content removes non-matching results and retains matching ones. & FR-5 \\
AT-6 & Ask a question & The answer is correct and every claim carries a timecode reference. & FR-6, NFR-5 \\
AT-7 & Unanswerable question & The system states that the footage does not support an answer instead of producing one. & NFR-5 \\
AT-8 & Play the moment & Selecting a result begins playback at that timecode without scrubbing, and several moments from one recording play consecutively as a reel. & FR-7, NFR-4 \\
AT-9 & Video-level review & Summary, chapters, events, entities and anomalies are produced and legible in the detail view. & FR-8 \\
AT-10 & Conversational retrieval & A multi-turn conversation selects the appropriate tools without the recording being named each turn, and is intact when resumed. & FR-9 \\
AT-11 & Index inspection & The stored record for any chunk is retrievable, and the timeline follows the playhead during playback. & FR-10 \\
AT-12 & Re-add and re-examine & The same file re-added is recognised as the same recording; re-examining at another depth preserves earlier results. & FR-11, NFR-10 \\
AT-13 & API discovery & A client constructed solely from the capability endpoint issues a valid search, and a newly registered analyzer appears without a client change. & FR-12, NFR-8 \\
AT-14 & Silent, cut-free footage & Fixed-camera footage with no speech still yields multiple chunks and usable descriptions. & NFR-6 \\
\bottomrule
\end{tabularx}
\end{table}

\subsection{Datasets}

\begin{table}[htbp]
\centering
\small
\begin{tabularx}{\textwidth}{@{}L{3.4cm}L{3.6cm}X@{}}
\toprule
\textbf{Dataset} & \textbf{Content} & \textbf{Property under test} \\
\midrule
Supermarket CCTV & 5 minutes, 1280$\times$720, fixed camera, no speech &
The principal fixture. A recording with neither cuts nor audio --- precisely the case
on which general-purpose tools fail. Used for chunking, scene, people, object and OCR
tests. \\
Multi-speaker clip & 60 seconds, several speakers &
The only fixture in which speaker change, silence and speaker-attributed transcription
all fire. Used for every audio path. \\
Roadway and incident clips & Short fixed-camera clips: licence plates, fire, suspected theft &
Exercise on-screen text recognition, novelty detection and event extraction on
distinct subject matter, and provide the moments used in clip-reel testing. \\
Extended set (planned) & Longer footage, multiple cameras, one shared subject &
Required for cross-recording person linking, which cannot be tested on footage sharing
no people. \\
\bottomrule
\end{tabularx}
\end{table}
